\chapter{PENDAHULUAN}
% Bab ini adalah bagian terpenting dari proposal. Dengan dukungan referensi
% yang kuat, tunjukkan URGENSI penelitian dan RESEARCH GAP yang hendak
% diselesaikan.

\section{Latar Belakang}
Latar belakang menjabarkan secara rinci konteks permasalahan yang akan diteliti. Dengan dukungan referensi yang kuat, pengusul harus mampu menunjukkan \textbf{urgensi} penelitian dan \textbf{\textit{research gap}} yang hendak diselesaikan. Tuliskan latar belakang secara mengalir (dari umum ke khusus) dengan menjawab pertanyaan-pertanyaan berikut:
\begin{enumerate}
  \item Apa fenomena atau permasalahan yang terjadi di dunia nyata atau di bidang keilmuan yang relevan? Dukung dengan data, statistik, atau fakta yang bersumber dari rujukan terpercaya.
  \item Mengapa permasalahan tersebut penting dan mendesak untuk diselesaikan (urgensi dari sisi akademik, industri, atau masyarakat)?
  \item Apa saja yang sudah dilakukan oleh peneliti terdahulu, dan apa keterbatasan pendekatan mereka?
  \item Celah penelitian (\textit{research gap}) apa yang belum terjawab oleh penelitian terdahulu?
  \item Pendekatan apa yang diusulkan untuk mengisi celah tersebut, dan apa kontribusi yang diharapkan?
\end{enumerate}
Akhiri latar belakang dengan pernyataan yang secara logis mengarah ke pertanyaan penelitian.

\section{Pertanyaan Penelitian}
Pertanyaan penelitian (\textit{research question}, RQ) merupakan inti persoalan yang hendak dijawab melalui penelitian dan diturunkan langsung dari \textit{research gap} pada latar belakang. Nyatakan dalam bentuk kalimat tanya yang spesifik dan dapat dijawab melalui rancangan penelitian pada Bab~\ref{chap:rancangan}, misalnya:
\begin{enumerate}[label=\textbf{RQ\arabic*}:, leftmargin=*]
  \item Bagaimana \textit{[metode/pendekatan]} dapat diterapkan untuk \textit{[menyelesaikan persoalan]} pada \textit{[objek/domain]}?
  \item Sejauh mana \textit{[metode/pendekatan]} meningkatkan \textit{[metrik kinerja]} dibandingkan dengan \textit{[pendekatan pembanding/baseline]}?
\end{enumerate}

\section{Tujuan Penelitian}
Tujuan penelitian mengungkapkan sasaran yang ingin dicapai. Setiap tujuan harus berkorespondensi satu-satu dengan pertanyaan penelitian dan dinyatakan menggunakan verba operasional yang terukur (misalnya: \textit{merancang, membangun, menganalisis, mengevaluasi}). Contoh:
\begin{enumerate}[label=\textbf{T\arabic*}:, leftmargin=*]
  \item Menerapkan \textit{[metode/pendekatan]} untuk \textit{[menyelesaikan persoalan]} pada \textit{[objek/domain]} (menjawab RQ1).
  \item Mengevaluasi kinerja \textit{[metode/pendekatan]} terhadap \textit{[baseline]} menggunakan metrik \textit{[metrik]} (menjawab RQ2).
\end{enumerate}

\section{Ruang Lingkup}
Ruang lingkup menegaskan batasan penelitian agar arah penelitian tetap fokus dan terukur. Tuliskan secara tegas dalam kalimat deklaratif, misalnya:
\begin{enumerate}
  \item Data yang digunakan adalah \textit{[nama/jenis dataset]} periode \textit{[rentang waktu]}.
  \item Metode yang digunakan dibatasi pada \textit{[nama metode/algoritma/framework]}.
  \item Penelitian dilakukan pada \textit{[unit/institusi/lingkungan penelitian]}.
  \item Evaluasi dilakukan menggunakan metrik \textit{[metrik]} dalam lingkungan \textit{[laboratorium/simulasi/lapangan terbatas]}.
\end{enumerate}
